\section{Planteamiento del Problema y Análisis del Contexto}
\label{sec:planteamiento}

\subsection{Descripción del Caso}
\label{subsec:descripcion-caso}

El Aeropuerto Internacional Jorge Chávez no distribuye a sus pasajeros de manera
uniforme. En determinadas franjas horarias y en puntos concretos del terminal se forman
concentraciones muy superiores a las del resto del recinto en ese mismo momento. Medir esa
combinación (cuántas personas, en qué zona, durante cuánto tiempo) importa porque las
decisiones que la atienden son de corto plazo y reversibles: habilitar un mostrador adicional,
redirigir el flujo hacia un filtro descargado, adelantar un relevo de personal o abrir una segunda
fila. Todas dependen de conocer dónde se está formando la concentración mientras se forma; un
recuento que llega al cierre de la jornada sirve para planificar el mes siguiente, pero no para
resolver la cola de esa mañana.

\subsection{Identificación de Variables y Restricciones}
\label{subsec:variables-restricciones}

\subsubsection{Variables Operativas de Salida}

\begin{itemize}
  \item \textbf{Variable de salida principal:} concentración de personas por zona (cuántas,
  dónde, durante cuánto tiempo).
  \item \textbf{Granularidad de salida:} admite dos niveles con consecuencias normativas
  distintas: un mapa de densidad agregado por zona, que no requiere identificar
  individuos, o una localización individual por puntos, que produce un dato
  espacio-temporal atribuible a una persona si se persiste y se encadena en el tiempo.
\end{itemize}

\subsubsection{Variables Técnicas del Pipeline}

Parámetros reales configurados en \texttt{proyecto\_lap\_prototipo/}:

\begin{itemize}
  \item \textbf{Detección:} umbral de confianza \texttt{umbral\_deteccion = 0.5}; resolución de
  entrada del detector configurable en 480, 640 o 960 px.
  \item \textbf{Seguimiento mono-cámara:} asociación húngara de alta confianza con
  \texttt{DISTANCIA\_MAX\_ALTA\_PX = 45.0}; filtro de Kalman de velocidad constante.
  \item \textbf{Fusión entre cámaras:} \texttt{umbral\_distancia\_fusion = 1.5 m};
  \texttt{ventana\_temporal\_s = 1.0 s} de recencia.
  \item \textbf{Filtrado geométrico:} \texttt{margen\_px = 15.0} de tolerancia en el borde del
  polígono de alcance.
  \item \textbf{Descriptor de apariencia:} sustracción de fondo MOG2 (\texttt{history=200},
  \texttt{varThreshold=25}, \texttt{detectShadows=False}).
  \item \textbf{Monitoreo integrado:} muestreo de 0.2 s por análisis; hasta 32 cámaras
  configurables por servidor (sin garantía de procesamiento simultáneo de todas).
\end{itemize}

\subsubsection{Variables de Desempeño y Evaluación}

Medidas, no estimadas:

\begin{itemize}
  \item Tiempo de inferencia de P2PNet en CPU: $\approx$3.5 s por cada 0.2 s de video de
  origen.
  \item Filtrado de identificadores: de 29 identificadores brutos a 10 personas reales en la
  corrida de validación.
  \item Fusiones cruzadas verificadas entre cámara A y B: 2, con distancia espacio-temporal
  promedio 0.76 y máxima 1.18 sobre un umbral de 1.5.
  \item Suite de pruebas del monitoreo integrado: 65 pruebas automatizadas, más compilación
  de TypeScript y Vite.
\end{itemize}

\subsection{Restricciones del Proyecto}
\label{subsec:restricciones}

\begin{itemize}
  \item \textbf{Restricciones Técnicas.} Esta entrega corresponde a un desarrollo y validación
  conceptual sobre datos públicos y grabaciones propias, sin acceso todavía a material
  del Aeropuerto Jorge Chávez; la literatura revisada documenta una brecha de dominio
  conocida entre esos conjuntos y un entorno de despliegue nuevo. El sistema debe
  operar en CPU, sin GPU dedicada, sobre cámaras fijas que no comparten resolución,
  altura ni ángulo de instalación.
  \item \textbf{Restricciones Éticas.} El sistema no puede derivar del video ningún rasgo que
  permita identificar a una persona, ni siquiera de forma indirecta (RNF-01). Un
  identificador persistente que permite reagrupar apariciones de un mismo individuo es
  disociación, no anonimización, y el dato disociado sigue siendo dato personal.
  \item \textbf{Restricciones Legales.} La Ley 29733 y su reglamento (D.S. 016-2024-JUS)
  exigen privacidad desde el diseño y por defecto. La Directiva 01-2020-JUS/DGTAIPD
  fija plazos de conservación de imágenes de videovigilancia entre 30 y 60 días,
  mientras que el D.S. 007-2020-IN (D.L. 1218 y Ley 30120, aplicable a
  establecimientos con aforo de 50 personas o más, umbral que un terminal aéreo
  supera con holgura) exige un mínimo de 45 días; ambos regímenes son compatibles
  solo dentro de una ventana estrecha, porque responden a finalidades distintas.
\end{itemize}

\subsection{Análisis Técnico, Ético, Social, Ambiental}
\label{subsec:analisis-teas}

\begin{itemize}
  \item \textbf{Análisis técnico.} La arquitectura separa deliberadamente cuántas personas hay
  (disponible incluso sin calibrar cámaras) de quién es quién entre cámaras (que exige
  calibración y sincronización declarada). La auditoría del prototipo muestra, sin
  embargo, que la implementación actual sí persiste la información necesaria para
  reconstruir la segunda, lo que reabre la pregunta de si esa capacidad debe exponerse
  en producción.

  \item \textbf{Análisis ético y legal.} El artículo 2 de la Constitución Política del Perú (1993)
  reconoce el derecho a que los servicios informáticos no afecten la intimidad y el
  derecho a la imagen propia. El sistema trata cuatro tipos de dato con régimen
  distinto: el fotograma (dato personal), el punto con coordenadas (no identifica por sí
  solo), el dato espacio-temporal (que encadenado reconstruye una trayectoria) y el
  dato biométrico propiamente dicho, que ni un mapa de densidad ni un conjunto de
  coordenadas de cabezas constituyen por sí mismos. El equipo detectó una
  discrepancia terminológica relevante: describir como personas anonimizadas el
  resultado de un proceso que asigna un identificador persistente es impreciso, porque
  la anonimización debe ser irreversible. De las diez medidas de privacidad desde el
  diseño que adopta el proyecto (M-01 a M-10), al cierre de esta versión se cumple
  M-02 (no persistencia del fotograma) pero no M-03 (salida agregada por zona sin
  coordenadas individuales persistidas): es el único hallazgo de incumplimiento que el
  equipo puede cerrar por sí solo, sin depender de LAP.

  \item \textbf{Análisis social.} Un sistema que no toma decisiones sobre personas individuales
  puede producir un daño distributivo si subcuenta sistemáticamente a un grupo:
  estatura, vestimenta, equipaje voluminoso o sillas de ruedas son factores
  identificables, aunque no cuantificados por ningún trabajo revisado. Existe además
  un riesgo de sustitución progresiva del criterio humano por el del sistema, si un
  operador aprueba mecánicamente las alertas automáticas.

  \item \textbf{Análisis ambiental.} El equipo no ha medido el consumo energético del sistema.
  Reutilizar modelos preentrenados sin reentrenar, procesar en el borde sin transmitir
  el fotograma completo, y no persistir imágenes de forma permanente reducen, al
  mismo tiempo, la exposición de datos personales y el cómputo necesario.
\end{itemize}
